Bei \thSO\ ist die Dichte des gesättigten Wasserdampfs $\rho_\text{s} = \rsSO$ (FoTaBe S.~192). Vor dem Aufguss enthält die Luft
\begin{align}
 m_1 &= \mDaF = \fiAO\cdot\rsS\cdot\VS = \mDa
\end{align}
Wasserdampf. Der Aufguss fügt $m_\text{W} = \mWaF = \rWa\cdot\VWa = \mWa$ hinzu. Bleibt die Temperatur gleich und kondensiert nichts, ist die relative Luftfeuchtigkeit anschliessend
\begin{align}
 \varphi_2 &= \frac{m_1+m_\text{W}}{\rho_\text{s}V} = \fiBF \\
   &= \fiAO + \frac{\mWa}{\rsS\cdot\VS} \\
   &= \fiB \approx \result{\fiBP{0}{2}}
\end{align}
Die relative Luftfeuchtigkeit verdoppelt sich etwa, bleibt aber sehr klein; trotzdem empfindet man den Aufguss als heiss, weil an der kühleren Haut Wasserdampf kondensiert.
